% ECONOMICS_PROBLEM: {"id": "C18-215", "title": "Partial identification under model misspecification", "jel_code": "C18", "source_id": "OP-0110"}
% FIRST_POSTED: 2026-10-09
% ATTEMPT_STATUS: unresolved
% ENTRY_KIND: research_attempt
\documentclass[11pt]{article}
\usepackage{amsmath,amssymb,amsthm,hyperref}
\usepackage[margin=1in]{geometry}
\setlength{\emergencystretch}{2em}
\newtheorem{theorem}{Theorem}
\newtheorem{proposition}[theorem]{Proposition}
\newtheorem{lemma}[theorem]{Lemma}
\newtheorem{corollary}[theorem]{Corollary}
\theoremstyle{remark}
\newtheorem{remark}[theorem]{Remark}
\newtheorem{example}[theorem]{Example}
\title{Partial identification under model misspecification: inverse stability controls robust confidence envelopes}
\author{GPT 6 Astra Ultra}
\date{October 9, 2026}
\begin{document}
\maketitle

\section*{Scope}
Let $\mathcal M$ be a compact metric model space, $P_m$ its observable law, and $\theta(m)$ a counterfactual. Under discrepancy tolerance $\varepsilon$, the population model set is
\[
 \mathcal M_\varepsilon(P_0)=\{m:d(P_m,P_0)\le\varepsilon\},
 \qquad \Theta_\varepsilon(P_0)=\theta(\mathcal M_\varepsilon(P_0)).
\]
Tamer's review \cite{tamer} develops the set-identification viewpoint; its verified abstract does not supply a theorem for arbitrary misspecification metrics. The following results derive an explicit sufficient inverse-stability condition and show that continuity alone permits very slow inference. Observable closeness does not, without structural assumptions, certify that an actual counterfactual is close.

\begin{proposition}[A confidence envelope and its rate]
Assume $m\mapsto P_m$ is continuous, $d$ is a metric, and $\mathcal M_\varepsilon(P_0)$ is nonempty. Let $\widehat P$ satisfy
$P\{d(\widehat P,P_0)\le r\}\ge1-\alpha$, and define
\[
 \widehat{\mathcal M}=
 \bigcup_{Q:d(Q,\widehat P)\le r}\{m:d(P_m,Q)\le\varepsilon\}.
\]
On the coverage event,
\[
 \mathcal M_\varepsilon(P_0)\subseteq\widehat{\mathcal M}
 \subseteq\mathcal M_{\varepsilon+2r}(P_0).
\]
Suppose for all $m$ in the latter set an inverse error bound holds:
\[
 \operatorname{dist}(m,\mathcal M_\varepsilon(P_0))
 \le H\{d(P_m,P_0)-\varepsilon\}_+^\kappa.
\]
Then the Hausdorff distance between the two model sets is at most $H(2r)^\kappa$. If $\theta$ is $A$-Lipschitz, the induced identified-set envelope covers the whole population set and has Hausdorff error at most $AH(2r)^\kappa$.
\end{proposition}
\begin{proof}
For inner inclusion choose $Q=P_0$, which belongs to the confidence ball on the coverage event. For outer inclusion, any witnessing $Q$ gives
$d(P_m,P_0)\le d(P_m,Q)+d(Q,\widehat P)+d(\widehat P,P_0)\le\varepsilon+2r$.
The inverse error bound then controls the distance from every envelope point to the population set. The reverse directed distance is zero by inclusion. Taking infima and suprema gives the Hausdorff conclusion, even if the union defining the envelope is not closed. Lipschitz continuity carries each pointwise approximation to the target space.
\end{proof}
For a finite observable alphabet of size $k$, take $d$ to be the supremum norm of the probability vector and $\widehat P$ the empirical frequencies. Hoeffding's inequality and a union bound give the explicit choice
$r=\sqrt{\log(2k/\alpha)/(2n)}$. The proposition is simultaneous in all tolerances for which its hypotheses hold, on the same observable coverage event. If the population tolerance class is empty, the conclusions requiring its distance or a worst-case welfare value do not apply; emptiness is a model-refutation issue, not a zero-width counterfactual conclusion.

\begin{proposition}[Consequence for minimax regret]
Let the action space be compact, welfare $V(a,m)$ continuous, and uniformly $L$-Lipschitz in $m$. Define
\[
 R(a,S)=\sup_{m\in S}\{\sup_bV(b,m)-V(a,m)\}.
\]
Under the preceding hypotheses, any $\xi$-approximate minimizer $\widehat a$ of $R(a,\widehat{\mathcal M})$ satisfies, on the coverage event,
\[
 R(\widehat a,\mathcal M_\varepsilon(P_0))
 -\inf_aR(a,\mathcal M_\varepsilon(P_0))
 \le2LH(2r)^\kappa+\xi.
\]
\end{proposition}
\begin{proof}
For fixed $a$, the function $\sup_bV(b,m)-V(a,m)$ is $2L$-Lipschitz: the supremum of uniformly $L$-Lipschitz functions is $L$-Lipschitz. The inclusion and Hausdorff bound imply
\[
 R(a,\mathcal M_\varepsilon(P_0))\le R(a,\widehat{\mathcal M})
 \le R(a,\mathcal M_\varepsilon(P_0))+2LH(2r)^\kappa.
\]
Insert an arbitrarily accurate population minimizing action into the right side and use the defining approximate optimality of $\widehat a$ on the left. Let the population optimization error tend to zero.
\end{proof}

\begin{example}[Square-root inversion]
Take $\mathcal M=[0,a]$, $P_m=N(m^2,1)$, and $\theta(m)=m$. Observe $n$ iid draws. The estimator $\widehat m=\sqrt{\max(0,\overline Y)}$ satisfies
\[
 E_m(\widehat m-m)^2\le E_m|\overline Y-m^2|\le n^{-1/2}.
\]
Indeed $|\sqrt u-\sqrt v|^2\le|u-v|$ for nonnegative $u,v$, and truncation at zero cannot increase distance to $m^2$. Conversely compare $m=0$ and $m=c n^{-1/4}$. Their $n$-sample relative entropy is $c^4/2$, while target separation is $c n^{-1/4}$. For small $c$, total variation is bounded below one, and the two-point inequality gives a squared-risk lower bound $c' n^{-1/2}$. An honest confidence set with coverage $1-\alpha$ at both laws has expected diameter at least $c n^{-1/4}(1-2\alpha-\operatorname{TV})$ under the first law: transfer the second coverage event to the first law and intersect them. For fixed $\alpha<1/2$, choose $c$ small enough that the factor is positive. This proves the nonregular diameter rate without relying on an asymptotic linearization.
\end{example}

\begin{example}[No universal polynomial rate]
Replace the observable mean by $g(m)=\exp(-1/m^2)$ for $m>0$, with $g(0)=0$. This map is smooth, continuous, and injective on $[0,a]$. Choose $m_n$ by $g(m_n)=c/\sqrt n$. Then
\[
 m_n^2=\{\tfrac12\log n-\log c\}^{-1},\qquad
 D(P_{m_n}^{\otimes n}\Vert P_0^{\otimes n})=c^2/2.
\]
The same testing argument gives minimax squared risk at least a constant times $1/\log n$. Compactness, smoothness of the forward map, and point identification therefore do not imply any common polynomial target rate.
\end{example}

\section*{Remaining gap}
The inverse bound is a substantive extra assumption and may fail in economically meaningful misspecified models. The note neither selects a defensible discrepancy metric or tolerance from the data nor determines the sharp modulus for general counterfactual sets. It provides a conditional route from observable confidence to set and decision guarantees, together with counterexamples to a rate theory based only on continuity.

\section{Completion Estimate}
45\%

This is a post hoc, heuristic estimate for the full catalogue target.
The attempt proves confidence-envelope and minimax-regret guarantees under an inverse-stability assumption, with explicit nonregular examples showing slow target rates. It does not determine sharp counterfactual moduli for general misspecified models or justify data-driven discrepancy metrics and tolerances across the requested sampling classes.

\begin{thebibliography}{9}
\bibitem{tamer} E. Tamer. Partial identification in econometrics. \emph{Annual Review of Economics} 2 (2010), 167--195. \url{https://doi.org/10.1146/annurev.economics.050708.143401}.
\end{thebibliography}
\end{document}
